\documentclass[12pt, a4paper]{article}
\usepackage[utf8]{inputenc}
\usepackage[T1]{fontenc}
\usepackage[czech]{babel}
\usepackage[top=25mm, bottom=25mm, left=25mm, right=25mm]{geometry}
\usepackage{lmodern}
\usepackage{microtype}
\usepackage{parskip}
\usepackage{enumitem}
\usepackage{hyperref}
\usepackage{fancyhdr}
\usepackage{titlesec}

\pagestyle{fancy}
\fancyhf{}
\fancyhead[L]{\nouppercase{\leftmark}}
\fancyhead[R]{\thepage}
\renewcommand{\headrulewidth}{0.4pt}

\title{\textbf{Průvodce projektem}\\[0.5em]
\large Systém pro správu skladu elektronických součástek\\[0.3em]
\normalsize Jak to funguje a jak to obhájit}
\author{Tomáš Simandl, Tomáš Szetei, Lukáš Kvapil}
\date{2026}

\begin{document}

\maketitle
\thispagestyle{empty}
\tableofcontents
\newpage

%% ── 1. O čem projekt je ──────────────────────────────────────────────────────
\section{O čem projekt je}

Navrhujeme informační systém pro sklad elektronických součástek. Sklad obsahuje různé součástky jako rezistory, kondenzátory nebo integrované obvody. Tyto součástky jsou fyzicky uloženy na konkrétních místech v regálech -- říkáme jim lokace. Součástky se naskladňují od dodavatelů, vydávají se na projekty a někdy se přesouvají mezi lokacemi.

Navíc lze součástky přiřazovat k projektům -- třeba projekt ``Robotická ruka'' potřebuje 50 kusů rezistoru a 10 kusů kondenzátoru. Systém hlídá kolik kusů je volných a kolik je rezervováno pro projekty.

Celý projekt je analytický návrh -- nekódujeme, ale kreslíme UML diagramy které přesně popisují jak by systém fungoval. Výsledkem jsou tři modely: objektový (class diagram), stavový (stavové diagramy) a model interakcí (use case, scénáře, sekvenční diagram, diagram aktivit).

%% ── 2. Základní pojmy ────────────────────────────────────────────────────────
\section{Základní pojmy}

\subsection{Třída a objekt}

Třída je šablona nebo vzor -- popisuje jak má objekt vypadat a co umí. Třída objektů popisuje skupinu objektů s podobnými vlastnostmi (atributy), společným chováním (operace) a stejnými vztahy k jiným objektům.

Jeden konkrétní objekt je \textbf{instance} třídy. Každá instance má své vlastní hodnoty atributů, ale sdílí společnou strukturu s ostatními instancemi té třídy.

Příklad: Třída \textbf{Soucastka} je šablona. Konkrétní ``Rezistor 10k ohm, 50 kusů, Dostupna'' je instance -- jeden konkrétní objekt vzniklý podle té šablony. Jiná instance je ``Kondenzátor 100uF, 20 kusů, Dostupna''.

\textbf{Důležité:} Dva objekty jsou různé přesto, že hodnoty jejich atributů mohou být identické -- každý objekt má vlastní identitu.

\subsection{Atribut}

Atribut popisuje datovou vlastnost objektu -- něco co objekt ``zná'' nebo ``má uložené''. Každý atribut má název a datový typ.

Datové typy:
\begin{itemize}
  \item \texttt{int} -- celé číslo. Například \texttt{mnozstvi: int} = počet kusů (50, 0, 200)
  \item \texttt{String} -- text. Například \texttt{nazev: String} = název součástky (``Rezistor 10k'')
  \item \texttt{boolean} -- pravda nebo nepravda (true/false)
  \item \texttt{Date} -- datum. Například \texttt{datum: Date} = kdy byl pohyb vytvořen
\end{itemize}

Třída \textbf{Soucastka} má atributy: \texttt{id: int}, \texttt{nazev: String}, \texttt{mnozstvi: int}, \texttt{rezervovaneMnozstvi: int}, \texttt{stav: String (Dostupna | Vydana | Odepsana)}.

\subsection{Metoda (operace)}

Metoda je implementace operace na třídu -- akce kterou objekt umí provést. Všechny objekty v třídě sdílejí stejné operace. Každá operace má cílový objekt jako implicitní argument a může mít další argumenty.

Příklady z projektu:
\begin{itemize}
  \item \texttt{zvysMnozstvi(pocet)} -- zvýší množství součástky na skladě o zadaný počet kusů. Volá se při příjmu zboží.
  \item \texttt{snizMnozstvi(pocet)} -- sníží množství. Volá se při výdeji.
  \item \texttt{rezervovat(pocet)} -- zablokuje kusy pro projekt. Součástka stále fyzicky je na skladě, ale systém ví že je rezervovaná.
  \item \texttt{zrusitRezervaci(pocet)} -- uvolní rezervaci. Kusy jsou opět volné.
  \item \texttt{overDostupnost(): boolean} -- zkontroluje zda je aspoň jeden kus volný (\texttt{mnozstvi - rezervovaneMnozstvi > 0}). Vrátí true nebo false.
\end{itemize}

Díky \textbf{polymorfismu} stejná operace může platit pro různé třídy a mít různé podoby. Operace \texttt{provestPohyb()} se chová jinak pro PrijemPohyb, VydejPohyb a PresunutiPohyb.

\subsection{Spojení a vazba (asociace)}

\textbf{Spojení} je fyzická nebo konceptuální souvislost mezi konkrétními objekty. Například: ``Skladník Jan Novák vytvořil příjem č. 42''.

\textbf{Vazba} (asociace) popisuje skupinu spojení se stejnou strukturou -- je to vztah mezi třídami. V diagramu ji vidíte jako čáru s popiskem. Vazba je svou podstatou dvousměrná.

Příklady z projektu:
\begin{itemize}
  \item \textbf{Uzivatel vytváří SkladovyPohyb} -- každý pohyb má svého autora. Bez uživatele by pohyb nevznikl.
  \item \textbf{SkladovyPohyb je součástí Soucastka} -- každý pohyb se týká konkrétní součástky.
  \item \textbf{SkladovyPohyb mění stav Sklad} -- pohyb změní stav skladu (zvýší nebo sníží zásoby).
  \item \textbf{Projekt přiřazuje Soucastka} -- projekt využívá určité součástky.
\end{itemize}

\subsection{Kardinalita (multiplicita)}

Kardinalita říká kolik objektů může být na každé straně vztahu. Píše se na konec čáry asociace.

\begin{itemize}
  \item \texttt{1} -- přesně jeden objekt
  \item \texttt{*} -- nula nebo více (libovolný počet)
  \item \texttt{0..1} -- žádný nebo jeden (volitelné)
  \item \texttt{1..*} -- jeden nebo více
\end{itemize}

V úvodních fázích nejdříve určíme objekty, třídy a vazby -- a až potom rozhodujeme o kardinalitě. Kardinalita není neměnná, upřesňuje se v průběhu analýzy.

Příklady z projektu:
\begin{itemize}
  \item \texttt{Uzivatel "1" --> "*" SkladovyPohyb} -- jeden uživatel může vytvořit libovolný počet pohybů. Jeden pohyb ale patří přesně jednomu uživateli.
  \item \texttt{Sklad "1" *-- "*" Lokace} -- jeden sklad má mnoho lokací. Jedna lokace patří přesně jednomu skladu.
  \item \texttt{Lokace "1" o-- "*" Soucastka} -- v jedné lokaci může být více součástek.
\end{itemize}

\subsection{Generalizace (dědičnost)}

Generalizace je vztah ``je speciálním typem''. Je to sdílení atributů a operací mezi třídami na hierarchickém principu. Nadtřída obsahuje obecnou informaci, kterou podtřídy upřesňují. Každá podtřída dědí všechny vlastnosti nadtřídy a přidává své vlastní.

V diagramu ho poznáte podle \textbf{prázdného trojúhelníku} na šipce směřující k nadřazené třídě.

V našem projektu: \texttt{SkladovyPohyb} je abstraktní nadtřída. Tři podtřídy:
\begin{itemize}
  \item \texttt{PrijemPohyb} -- dědí vše + přidává \texttt{dodavatel: String}
  \item \texttt{VydejPohyb} -- dědí vše + přidává \texttt{duvod: String}
  \item \texttt{PresunutiPohyb} -- dědí vše + přidává \texttt{zdrojLokaceId: int} a \texttt{cilLokaceId: int}
\end{itemize}

\textbf{Abstraktní třída} je třída ze které nelze vytvořit instanci. \texttt{SkladovyPohyb} je abstraktní -- v reálu nikdy nevznikne ``obecný pohyb'', vždy je to konkrétně příjem, výdej nebo přesun.

\textbf{Proč generalizaci používáme?} Příjem, výdej i přesun sdílejí atributy \texttt{id}, \texttt{datum}, \texttt{stav}, \texttt{mnozstvi} a metody \texttt{vytvoritPohyb()}, \texttt{overitData()}, \texttt{provestPohyb()}. Bez dědičnosti bychom je museli opisovat třikrát. Generalizace je jednou ze základních výhod objektového přístupu -- velmi zmenšuje opakování v návrhu.

\subsection{Agregace -- části přežijí bez celku}

Agregace je zvláštní typ vazby -- vztah ``součást-celek''. Část může existovat i bez nadřazeného celku. Symbolem agregace je \textbf{prázdný kosočtverec} (diamant) na straně celku. Agregace je tranzitivní a antisymetrická.

V našem projektu: \texttt{Lokace o-- Soucastka} -- lokace ukládá součástky, ale součástka může existovat bez lokace. Například: přijde zásilka rezistorů, jsou zaevidovány v systému, ale skladník je ještě fyzicky neodnesl na místo v regálu. Součástka existuje, lokaci nemá -- a to je v pořádku.

\textbf{Co se stane když lokace zanikne?} Součástky přijdou o přiřazení k lokaci, ale stále existují v systému.

\subsection{Kompozice -- části zaniknou s celkem}

Kompozice je zvláštní případ agregace pro který platí dvě omezení: součást může patřit pouze jednomu celku a zánikem celku zaniká i součást. Symbolem je \textbf{plný (vyplněný) kosočtverec}.

V našem projektu: \texttt{Sklad *-- Lokace} -- sklad se skládá z lokací. Lokace ``Regál A, Police 3'' nedává smysl bez konkrétního skladu. Pokud by sklad zanikl, zaniknou i všechny jeho lokace.

\textbf{Rozdíl agregace vs kompozice:} U agregace části přežijí (součástka bez lokace). U kompozice části zaniknou (lokace bez skladu).

\subsection{Kvalifikovaná vazba}

Kvalifikovaná vazba dává do vztahu dvě třídy a \textbf{kvalifikátor} -- zvláštní atribut, který účinně zmenšuje násobnost vazby z 1:N na 1:1. Kvalifikátor rozlišuje mezi objekty na konci ``mnoho'' vazby. Kreslí se jako \textbf{rámeček} na tom konci vazby, který kvalifikuje.

V našem projektu: \texttt{Sklad "[id: int]" --> Soucastka} -- sklad přistupuje ke konkrétní součástce pomocí jejího ID. Metoda \texttt{najdiSoucastku(id): Soucastka} vrátí přesně tu jednu součástku se zadaným ID. Bez kvalifikátoru bychom přistupovali k celé kolekci součástek najednou.

Kvalifikace zlepšuje sémantickou přesnost a názornost vazby.

\subsection{Vazební (asociační) třída}

Vazební třída je vazbou a současně třídou -- je to třída která patří vazbě, ne jednomu z propojených objektů. Atribut spojení je vlastností vazby, ne objektu. Má opodstatnění hlavně ve vazbách mnoho-mnoho. V diagramu je napojena \textbf{čárkovanou linkou} přímo na čáru asociace.

V našem projektu: \texttt{ProjektSoucastka} s atributem \texttt{mnozstvi: int}. Projekt ``Robotická ruka'' potřebuje 50 kusů rezistoru R1. Jiný projekt potřebuje 10 kusů stejného rezistoru. Toto množství nepatří ani Projektu (ten má více součástek) ani Součástce (ta je ve více projektech) -- patří konkrétní kombinaci projekt + součástka.

\textbf{Může vazební třída existovat bez obou propojených objektů?} Ne. Vznikne jen tehdy když existuje konkrétní projekt a konkrétní součástka a jsou svázány. Zanikne-li projekt nebo součástka, zanikne i příslušný záznam.

\subsection{Polymorfismus}

Polymorfismus znamená že stejná operace se může chovat odlišně v různých třídách. Specifická implementace operace určitou třídou se nazývá metoda. Operace může být implementována více než jednou metodou.

V našem projektu: voláme \texttt{provestPohyb()} na pohybu:
\begin{itemize}
  \item \texttt{PrijemPohyb.provestPohyb()} -- zavolá \texttt{zvysMnozstvi()} a zvýší množství
  \item \texttt{VydejPohyb.provestPohyb()} -- zavolá \texttt{snizMnozstvi()} a sníží množství
  \item \texttt{PresunutiPohyb.provestPohyb()} -- zavolá \texttt{zmenLokaci()} a přesune součástku
\end{itemize}

Stejné volání, tři různé výsledky. Kód který pohyb volá nemusí vědět o jaký typ jde.

%% ── 3. Stavové diagramy ──────────────────────────────────────────────────────
\section{Stavové diagramy -- Dynamický model}

Stavový diagram ukazuje vztah událostí a stavů. Když nastane událost, další stav závisí na této události a současném stavu. Stavový diagram popisuje chování jediné třídy objektů.

\textbf{Notace:}
\begin{itemize}
  \item Stav: zaoblený obdélník s názvem
  \item Přechod: šipka, nápis je název události
  \item Výchozí stav: plný kroužek
  \item Koncový stav: plný kroužek v kroužku (``volské oko'')
\end{itemize}

\subsection{Události}

Událost je něco co se stane v časovém \textbf{okamžiku}. Událost nemá žádné časové trvání. Například ``Pohyb odeslán'', ``Chyba ověření'', ``Vydána ze skladu''.

Typy událostí:
\begin{itemize}
  \item \textbf{Signálová událost} -- explicitní přenos informace od jednoho objektu k druhému. Například ``Pohyb odeslán'' je signál od uživatele k systému.
  \item \textbf{Změnová událost} -- nastane když se hodnota booleovského výrazu změní z false na true. Například \texttt{when(mnozstvi < minimum)}.
  \item \textbf{Časová událost} -- nastane po uplynutí časového intervalu. Například \texttt{after(10 sekund)}.
\end{itemize}

\subsection{Podmínky}

Podmínka je booleovský výraz hodnot objektu. Na rozdíl od události (která nemá trvání) podmínka platí v \textbf{časovém intervalu}. Podmínky se používají jako hlídače na přechodech -- píší se v hranatých závorkách za názvem události.

Přechod se odpálí pouze pokud je podmínka splněna \textit{v okamžiku příchodu události}. Pokud je podmínka splněna až později, přechod se neodpálí.

Příklady z projektu: \texttt{[typ == příjem]}, \texttt{[mnozstvi > 0]}, \texttt{[neplatná data]}.

\subsection{Aktivity -- entry, do, exit}

Popis chování objektu musí určit co objekt dělá jako odezvu na událost. Aktivity se provádějí na přechodu nebo na začátku, konci či uvnitř stavu. Notace: ``/ jméno aktivity'' za názvem události.

Pořadí aktivit ve stavu:
\begin{enumerate}
  \item Aktivity na vstupním přechodu
  \item \textbf{entry/} -- vstupní aktivity (provedou se při vstupu do stavu)
  \item \textbf{do/} -- do-aktivity (probíhají po celou dobu pobytu ve stavu)
  \item \textbf{exit/} -- výstupní aktivity (provedou se při odchodu ze stavu)
  \item Aktivity na výstupním přechodu
\end{enumerate}

\textbf{Do-aktivita} je operace která má \textit{časové trvání}. Obsahuje spojité operace (trvají bez omezení) nebo sekvenční operace (po době samovolně skončí). Do-aktivitu může přerušit událost přijatá v průběhu jejího provádění -- ale výstupní aktivita se ještě provede.

\textbf{Důležité:} Okamžité akce (bez trvání) jako \texttt{zvysMnozstvi()} nepatří do \texttt{do/} -- patří do \texttt{entry/} nebo jako vnitřní přechod.

Příklad z projektu (stav Validovan):
\begin{itemize}
  \item \texttt{entry / overitData()} -- při vstupu se hned ověří data
  \item \texttt{do / kontrolovat data} -- průběžně probíhá kontrola
  \item \texttt{exit / potvrdit výsledek} -- před odchodem se potvrdí výsledek
\end{itemize}

\subsection{Vnitřní přechod}

Vnitřní přechod je reakce na událost která \textbf{nezmění stav} -- objekt zůstane ve stejném stavu ale provede okamžitou akci. Zapisuje se ve tvaru ``událost / aktivita'' uvnitř symbolu stavu.

Příklad z projektu (stav Dostupna součástky):
\begin{itemize}
  \item \texttt{naskladněno / zvysMnozstvi(pocet)} -- přijde zásilka, zvýší se množství, ale součástka zůstane Dostupna
  \item \texttt{rezervace přijata / rezervovat(pocet)} -- část kusů se zablokuje pro projekt, stav se nemění
  \item \texttt{rezervace zrušena / zrusitRezervaci(pocet)} -- rezervace se uvolní, stav se nemění
\end{itemize}

\subsection{Zakončovací přechod}

Zakončovací přechod se odpálí automaticky po skončení do-aktivity stavu -- bez čekání na událost. Přechod \textbf{nemá název události}. Pokud není splněna podmínka na zakončovacím přechodu, nastane ``zamrznutí'' ve stavu.

Příklad z projektu: stav \texttt{Proveden} provede \texttt{entry/provestPohyb()}, průběžné akce a \texttt{exit/uložit záznam} -- a pak automaticky přejde do koncového stavu bez čekání.

%% ── 4. Model interakcí ───────────────────────────────────────────────────────
\section{Model interakcí}

Model interakcí říká jak na sebe objekty navzájem působí -- je to celostní pohled přes mnoho objektů. Stavový model popisuje chování objektů jednotlivě, model interakcí popisuje jejich spolupráci. K úplnému popisu chování potřebujeme oba modely.

Interakce modelujeme na třech úrovních abstrakce:
\begin{itemize}
  \item \textbf{Use Case} -- vrcholová úroveň, interakce systému s okolím
  \item \textbf{Sekvenční diagram} -- výměna zpráv v čase v rámci množiny objektů
  \item \textbf{Diagram aktivit} -- tok řízení mezi výpočetními kroky (nejdetailnější)
\end{itemize}

\subsection{Use Case}

Aktor je vnější uživatel systému, který se systémem přímo komunikuje ale není součástí systému. Aktor má jediný a dobře definovaný účel. Aktor nemusí být pouze osoba -- může to být i jiný systém nebo zařízení.

Use Case je ucelená část funkčnosti, kterou systém poskytuje při komunikaci s aktorem. Každý use case musí poskytovat uživatelům určitou souhrnnou službu -- nesmí být definován příliš úzce.

\textbf{Vztahy v Use Case:}
\begin{itemize}
  \item \textbf{Include} -- jeden use case zahrnuje jiný use case. Zahrnutí je \textit{povinné} -- bez zahrnutého UC nemůže původní existovat. Notace: přerušovaná šipka \texttt{<<include>>}. Příklad: Vytvořit pohyb vždy zahrnuje Ověřit data pohybu.
  \item \textbf{Extend} -- přidá use case další \textit{volitelnou} funkčnost. Původní UC dává smysl i bez přidaného. Notace: přerušovaná šipka \texttt{<<extend>>}.
  \item \textbf{Generalizace} -- vyjadřuje specifické varianty use case. Rodičovský UC je obvykle abstraktní. Notace: prázdný trojúhelník (stejná jako v class diagramu).
\end{itemize}

\textbf{Proč v našem projektu generalizace místo extend?} Příjem, výdej a přesun jsou specializace pohybu -- ne volitelná rozšíření. Zavolání ``Vytvořit pohyb'' bez konkrétního typu nedává smysl. Generalizace navíc ladí s class diagramem kde stejný vztah používáme mezi \texttt{SkladovyPohyb} a podtřídami.

\subsection{Sekvenční diagram}

Sekvenční diagram ukazuje účastníky interakce a zprávy mezi nimi. Každý aktor je reprezentován svislou čarou (čára života) a každá zpráva je horizontální šipka od odesilatele k příjemci. Čas postupuje shora dolů.

\textbf{Typy šipek:}
\begin{itemize}
  \item \textbf{Plná šipka} (\texttt{->}) -- synchronní zpráva. Odesílatel čeká na odpověď.
  \item \textbf{Přerušovaná šipka} (\texttt{-->}) -- odpověď nebo návratová hodnota.
\end{itemize}

\textbf{Activation boxy} (obdélníky na čárách života) ukazují kdy je objekt aktivní -- kdy zpracovává zprávu.

\textbf{Alt blok} zobrazuje alternativní toky. Podmínka v hranatých závorkách určuje která větev se provede.

Pro každý use case je vhodné udělat alespoň jeden sekvenční diagram -- pro hlavní tok a pro každý chybový případ.

\subsection{Diagram aktivit}

Diagram aktivit ukazuje tok řízení -- kroky, větvení a slučování. Je podobný vývojovému diagramu, ale navíc může modelovat i konkurenční toky.

\textbf{Notace:}
\begin{itemize}
  \item Aktivita: zaoblený obdélník s názvem
  \item Rozhodovací uzel: diamant (kosočtverec)
  \item Podmínka: v hranatých závorkách \texttt{[ano]}, \texttt{[ne]}
  \item Začátek: plné kolečko
  \item Konec: plné kolečko v kroužku
\end{itemize}

Podmínky se testují po dokončení aktivity. Je vhodné používat podmínku \texttt{[else]} aby bylo zaručeno pokračování.

%% ── 5. Jak modely navazují ───────────────────────────────────────────────────
\section{Jak modely na sebe navazují}

Toto je klíčová otázka u zkoušky. Všechny tři modely jsou potřebné -- poskytují pohled na systém ze tří různých hledisek:

\begin{itemize}
  \item \textbf{Objektový model (class diagram)} -- statická struktura. Definuje třídy, vazby a zobecnění. Je to základ -- všechny ostatní diagramy z něj vycházejí. Metody které vidíme ve stavových a sekvenčních diagramech jsou přesně ty metody definované v class diagramu.
  \item \textbf{Stavový model} -- změny vlastností objektů v čase. Stavové diagramy pro \texttt{SkladovyPohyb} a \texttt{Soucastka} ukazují životní cykly těchto tříd z class diagramu.
  \item \textbf{Model interakcí} -- spolupráce objektů. Use Case ukazuje co systém umí. Sekvenční diagram a diagram aktivit popisují stejnou operaci (vytvoření pohybu) z různých pohledů -- sekvenční zdůrazňuje kdo s kým komunikuje, aktivit zdůrazňuje tok řízení.
\end{itemize}

Generalizace je vztah ``NEBO'' -- vyskytuje se ve všech třech modelech. Všechny superklasifikátory jsou abstraktní, listové klasifikátory jsou konkrétní.

Agregace je vztah ``A'' s vlastností tranzitivity a antisymetrie. Má dvě podoby: obecnou agregaci (složky mohou existovat nezávisle na celku) a kompozici (složka náleží jednomu celku a má stejnou dobu existence).

%% ── 6. Typické otázky u zkoušky ──────────────────────────────────────────────
\section{Typické otázky u zkoušky}

\textbf{„Jak spolu vaše diagramy souvisí?"}\\
Objektový model definuje třídy -- ty se pak objevují ve stavových diagramech (životní cykly), sekvenčním diagramu (volání metod) i diagramu aktivit. Use Case generalizace ladí s generalizací v class diagramu. Scénáře a diagramy popisují stejné operace z různých pohledů.

\textbf{„Proč je SkladovyPohyb abstraktní?"}\\
Protože obecný pohyb v reálu neexistuje -- vždy je to příjem, výdej nebo přesun. Abstraktní třída sdílí společné atributy a deklaruje metody které každá podtřída musí implementovat po svém. Z abstraktní třídy nelze vytvořit instanci.

\textbf{„Co je kvalifikovaná vazba a proč ji máte?"}\\
Kvalifikátor je zvláštní atribut který zmenšuje násobnost vazby z 1:N na 1:1. Sklad přistupuje ke konkrétní součástce pomocí jejího ID -- přímý přístup místo prohledávání celé kolekce. Kreslí se jako rámeček na konci vazby.

\textbf{„Proč máte vazební třídu ProjektSoucastka?"}\\
Atribut \texttt{mnozstvi} nepatří ani Projektu ani Součástce -- patří konkrétní vazbě mezi nimi. Projekt ``Robotická ruka'' potřebuje 50 kusů rezistoru, jiný projekt potřebuje 10 kusů stejného rezistoru. Takový atribut je vlastností spojení a nelze ho přiřadit k žádnému z objektů.

\textbf{„Jaký je rozdíl mezi agregací a kompozicí?"}\\
Agregace -- části mohou existovat bez celku (součástka může být bez lokace při příjmu). Kompozice -- části zaniknou spolu s celkem (lokace bez skladu nedává smysl). Oba jsou vztahy celek-část, liší se mírou závislosti.

\textbf{„Co je polymorfismus?"}\\
Stejná operace \texttt{provestPohyb()} se chová různě podle podtřídy -- příjem zvýší množství, výdej sníží, přesun změní lokaci. Polymorfismus nahrazuje větvení dle typu -- kód nemusí vědět o jaký typ pohybu jde.

\textbf{„Co znamená entry/do/exit?"}\\
Entry se provede při vstupu do stavu. Do probíhá po celou dobu pobytu (musí mít časové trvání). Exit se provede při odchodu. Pořadí: vstupní přechod, entry, do, exit, výstupní přechod.

\textbf{„Proč zakončovací přechody bez názvu?"}\\
Stavy Proveden a Zamitnut jen provedou akce a automaticky skončí -- nečekají na žádnou událost. Zakončovací přechod se odpálí automaticky po skončení do-aktivity.

\textbf{„Co je vnitřní přechod a kdy ho použijeme?"}\\
Reakce na událost která nezmění stav. Součástka zůstane Dostupna ale při události ``naskladněno'' se zavolá \texttt{zvysMnozstvi()} a zvýší se množství. Vnitřní přechod se používá pro okamžité akce bez trvání, které nemění stav objektu.

\textbf{„Proč generalizace v Use Case místo extend?"}\\
Include je povinné zahrnutí, extend je volitelné rozšíření. Příjem/výdej/přesun jsou specializace pohybu -- ne volitelná rozšíření. Rodičovský UC ``Vytvořit pohyb'' by bez nich nedával smysl. Generalizace navíc ladí s class diagramem.

\textbf{„Jaký je rozdíl mezi sekvenčním diagramem a diagramem aktivit?"}\\
Sekvenční diagram -- kdo s kým komunikuje (objekty a zprávy mezi nimi). Diagram aktivit -- tok řízení (kroky a větvení). Oba popisují vytvoření pohybu, ale z různých pohledů. Pro každý use case je vhodný alespoň jeden sekvenční diagram -- pro hlavní tok a pro každou chybu.

\textbf{„Co je událost a liší se od podmínky?"}\\
Událost nastane v okamžiku a nemá trvání (``Pohyb odeslán''). Podmínka platí v časovém intervalu -- testuje se jednou v okamžiku příchodu události. Podmínka se píše v hranatých závorkách za názvem události.

\textbf{„Proč nesmíte dát do-aktivitu pro okamžitou akci jako zvysMnozstvi()?"}\\
Do-aktivita musí mít časové trvání -- probíhá po celou dobu pobytu ve stavu. Akce \texttt{zvysMnozstvi()} je okamžitá -- proběhne a je hotová. Taková akce patří do \texttt{entry/} nebo jako vnitřní přechod. Pokud bychom ji dali do \texttt{do/}, systém by ji prováděl opakovaně po celou dobu pobytu ve stavu -- což nedává smysl.

\end{document}
